Government endpoint fleets face more disclosed vulnerabilities than they can
remediate within operational capacity, and prioritization decisions are often made
without traceable, reviewable evidence. We present an open, reproducible benchmark
framework for context-aware vulnerability-host \textit{pair} prioritization that
integrates exploit-intelligence (EPSS, KEV, public proof-of-concept signals),
asset-criticality, and endpoint-exposure features, and that couples ranking with a
capacity-constrained scheduling simulation and an append-only, hash-chained audit
log of every decision. We evaluate the framework on a deterministic,
public-sector-shaped synthetic fleet using a frozen 30-seed artifact spanning 13
prioritization strategies and a simulated expected-exploited-host-days (EHD)
operational metric, with strict output inspection and a content-addressed
freeze/verify protocol. The evaluation checks artifact integrity end to end (4,290
metric rows with no missing or non-finite values; 390 audit logs that all verify;
feasible scheduling within capacity) and computes operational and ranking metrics
reproducibly from the frozen outputs. Under the current synthetic fixtures and
uncalibrated (placeholder) weights, the strategies are statistically
indistinguishable on EHD: the proposed context-aware model neither beats the EPSS
baseline nor a random ordering, and all inter-strategy differences fall within
seed-to-seed variation. These results expose capacity-driven and base-rate
trade-offs and, rather than establishing real-world superiority, provide a
falsifiable, audit-evidence-producing benchmark on which calibrated and externally
validated prioritization studies can build. *(All numeric values to be reconfirmed
against the released artifact before camera-ready.)*

management; reproducible benchmark; auditability; public-sector security operations.
